% ===================================================================
%  Fig. 1 — Arquitectura logica propuesta
%  Entregable: short-paper-es
%  Figura vectorial en TikZ: se edita como texto y compila con el paper.
%  Preambulo requerido:
%    \usepackage{tikz}
%    \usetikzlibrary{arrows.meta,positioning,fit,backgrounds,shapes.geometric}
%  Va dentro de un entorno figure* (ancho de las dos columnas).
% ===================================================================
\begin{tikzpicture}[
  font=\footnotesize,
  % Estilos base
  banda/.style   = {draw=black!40, rounded corners=4pt, line width=0.8pt},
  caja/.style    = {draw=black!60, fill=white, rounded corners=3pt,
                    align=center, minimum height=10mm, inner sep=4pt,
                    line width=0.6pt},
  cliente/.style = {caja, fill=orange!15, text width=25mm},
  servicio/.style= {caja, fill=blue!12, text width=28mm},
  persist/.style = {caja, fill=cyan!10, text width=30mm},
  api/.style     = {caja, fill=blue!20, text width=35mm, minimum height=14mm},
  etiq/.style    = {font=\scriptsize\bfseries, black!70},
  fl/.style      = {-{Latex[length=2.5mm]}, line width=0.7pt, black!70},
  flInv/.style   = {Latex[length=2.5mm]-, line width=0.7pt, black!70},
  flBid/.style   = {Latex[length=2.5mm]-Latex[length=2.5mm], line width=0.7pt, black!70},
  nota/.style    = {font=\tiny\itshape, black!50, align=center, fill=white,
                    inner sep=2pt, rounded corners=2pt},
]

% =============== CAPA 1: CLIENTE (arriba) ===============
\begin{scope}[on background layer]
  \node[banda, fill=gray!5, minimum width=160mm, minimum height=40mm,
        rounded corners=5pt] at (8,8) {};
\end{scope}
\node[etiq] at (8,9.8) {\textbf{Cliente}};

% Nodos del cliente
\node[cliente, fill=orange!25, text width=28mm, minimum height=15mm] (ficha) at (2.5,8)
  {\textbf{Ficha de}\\[-1pt]\textbf{Producto}\\[-3pt]
   \tiny\textbf{DOM}};

\node[cliente, fill=green!15, text width=30mm, minimum height=15mm] (extn) at (8,8)
  {\textbf{Extensión de}\\[-1pt]\textbf{navegador}};

\node[cliente, fill=cyan!15, text width=28mm, minimum height=15mm] (app) at (13.5,8)
  {\textbf{Aplicación}\\[-1pt]\textbf{web}};

% Flechas cliente
\draw[fl] (ficha) -- node[nota, above, yshift=1mm] {Envía información} (extn);
\draw[flBid] (extn) -- node[nota, above, yshift=1mm] {Solicita autenticación} (app);

% =============== BANDA DE CONEXIÓN ===============
\draw[fl] (extn.south) -- ++(0,-1.2) -| node[nota, above, pos=0.2] {Envía información extraída de la ficha de producto} (8,5.8);

% =============== CAPA 2: MOTOR DE RIESGO (medio) ===============
\begin{scope}[on background layer]
  \node[banda, fill=gray!5, minimum width=160mm, minimum height=60mm,
        rounded corners=5pt] at (8,3.5) {};
\end{scope}
\node[etiq] at (8,6.2) {\textbf{Motor de Riesgo}};

% API centralizadora
\node[api, fill=blue!25, minimum height=16mm] (api) at (8,4.5)
  {\textbf{API centralizadora de estimación}\\[-1pt]
   \textbf{del motor de riesgo}};

% Servicios del motor de riesgo
\node[servicio, fill=blue!15, text width=25mm, minimum height=12mm] (dom) at (1.5,2)
  {\textbf{Analizador de veracidad}\\[-1pt]\textbf{de dominio}};

\node[servicio, fill=blue!15, text width=25mm, minimum height=12mm] (rep) at (5.5,2)
  {\textbf{Analizador de}\\[-1pt]\textbf{reputación}};

\node[servicio, fill=blue!15, text width=28mm, minimum height=12mm] (rf) at (10.5,2)
  {\textbf{API de estimación del nivel}\\[-1pt]\textbf{de probabilidad de riesgo}};

\node[servicio, fill=blue!15, text width=28mm, minimum height=12mm] (bert) at (14.5,2)
  {\textbf{API de Detección de}\\[-1pt]\textbf{reseñas falsas}};

% Flechas de API a servicios
\draw[fl] (api.south) -- ++(0,-0.5) -| node[nota, above, pos=0.3] {Solicita análisis de veracidad de dominio} (dom.north);
\draw[fl] (api.south) -- ++(0,-0.5) -| node[nota, above, pos=0.3] {Solicita análisis de reputación del vendedor} (rep.north);
\draw[fl] (api.south) -- ++(0,-0.5) -| node[nota, above, pos=0.3] {Solicita análisis del nivel de probabilidad de riesgo} (rf.north);
\draw[fl] (api.south) -- ++(0,-0.5) -| node[nota, above, pos=0.3] {Solicita porcentaje de reseñas falsas} (bert.north);

% =============== BANDA DE CONEXIÓN A PERSISTENCIA ===============
\draw[fl] (api.south) -- ++(0,-0.8) -| node[nota, above, pos=0.2] {Almacena información de consultas realizadas} (8,-0.8);

% =============== CAPA 3: PERSISTENCIA (abajo) ===============
\begin{scope}[on background layer]
  \node[banda, fill=gray!5, minimum width=160mm, minimum height=35mm,
        rounded corners=5pt] at (8,-2.5) {};
\end{scope}
\node[etiq] at (8,-0.8) {\textbf{Persistencia}};

% Base de datos
\node[persist, fill=cyan!15, text width=35mm, minimum height=18mm] (bd) at (8,-3)
  {\textbf{Base de Datos}\\[-3pt]
   \tiny Consultas · características · veredictos\\[-1pt]
   \tiny (trazabilidad y reentrenamiento)};

% Flecha a base de datos
\draw[fl] (8,-1.5) -- (bd.north);

\end{tikzpicture}
